\documentclass[12pt]{article}
\usepackage[utf8]{inputenc}
\usepackage[T1]{fontenc}
\usepackage{hyperref}
\title{Galatian Language Resources for K-12 Curriculum Development}
\author{Compiled by Claude Sonnet 5.5 (claude-sonnet-5-5) for Andr\'as Kornai}
\date{}
\begin{document}
\maketitle

\section{Introduction}
Galatian (ISO 639-3 xga; Glottocode gala1252) was a Celtic language of central Anatolia, probably close to Gaulish. It was used from the 3rd century BC to at least the 4th century AD. No Galatian texts survive. The evidence is a few glosses in classical authors and personal and place names, about 120 words in all. No revitalization programme or school materials were found.

\section{Primary Educational Resources}
None found.

\section{Linguistic and Reference Materials}
\begin{itemize}
\item Wikipedia article on Galatian: \url{https://en.wikipedia.org/wiki/Galatian_language}
\item Wikipedia article on the Continental Celtic languages: \url{https://en.wikipedia.org/wiki/Continental_Celtic_languages}
\item Bryn Mawr Classical Review 2002.12.19, a review noting that only personal names survive: \url{https://bmcr.brynmawr.edu/?p=22725}
\item University of British Columbia repository item: \url{https://circle.ubc.ca/handle/2429/21055}. I did not confirm its content.
\item Overview page on the Galatians as Celts in Anatolia: \url{https://folkdancefootnotes.org/culture/ethnicity-history-geography/galatians-celts-in-anatolia/}
\end{itemize}

\section{Religious and Community Resources}
None found for the language itself.

\section{Digital and Audio Resources}
None found.

\section{Curriculum Development Context}
Too little Galatian survives to support a language curriculum. Any school use would be a history or classics unit on the Galatians, with Gaulish and other Celtic languages as comparative context. Material beyond the names and glosses would be reconstruction.
\end{document}